\documentclass[10pt]{article}
\usepackage[T1]{fontenc}
\usepackage[utf8]{inputenc}
\usepackage[french]{babel}
\usepackage{amsmath,amssymb,mathrsfs,mathpazo}
\usepackage{geometry}
\geometry{a4paper,margin=2cm}
\usepackage{graphicx,array,tabularx,tabularray,colortbl,multirow,multicol,booktabs,enumitem,pifont,ulem,fancybox,diagbox,cancel}
\usepackage{pstricks,pstricks-add,pst-plot,pst-tree,pst-node,pst-text,pst-func,pst-math,pst-bspline,pst-3dplot,pst-eucl,pst-solides3d}
\let\pscalculate\undefined
\usepackage{pst-bezier}
\usepackage{tikz,pgfplots,tkz-tab}
\usetikzlibrary{arrows,arrows.meta,patterns,calc,positioning,trees,decorations.markings,shapes.geometric,shapes.misc,shapes.arrows,intersections,decorations.pathmorphing,decorations.pathreplacing,backgrounds,shadows,babel}
\pgfplotsset{compat=1.18}
\usepgfplotslibrary{fillbetween}
\usepackage[np]{numprint}
\usepackage{listings,makecell,pas-tableur,twoopt}
\let\barre\undefined
\usepackage[pstricks]{tablvar}
\definecolor{darkgreen}{RGB}{0,100,0}
\providecommand{\red}{\color{red}}
\providecommand{\blue}{\color{blue}}
\providecommand{\green}{\color{green}}
\providecommand{\black}{\color{black}}
\providecommand{\white}{\color{white}}
\providecommand{\gray}{\color{gray}}
\providecommand{\cyan}{\color{cyan}}
\providecommand{\magenta}{\color{magenta}}
\providecommand{\yellow}{\color{yellow}}
\providecommand{\violet}{\color{violet}}
\providecommand{\orange}{\color{orange}}
\providecommand{\euro}{\texteuro}
\providecommand{\EUR}{\texteuro}
\providecommand{\ieme}{\textsuperscript{e}}
\providecommand{\degres}{\ensuremath{^\circ}}
\providecommand{\degree}{\ensuremath{^\circ}}
\providecommand{\vv}[1]{\overrightarrow{#1}}
\providecommand{\Coord}[1]{#1}
\providecommand{\decosix}{}
\providecommand{\starredbullet}{\ensuremath{\bullet}}
\pagestyle{empty}
\setlength{\parindent}{0pt}
\usepackage{adjustbox}\DeclareUnicodeCharacter{2236}{:}\DeclareUnicodeCharacter{2190}{\ensuremath{\leftarrow}}
\providecommand{\R}{\mathbb{R}}
\providecommand{\N}{\mathbb{N}}
\providecommand{\Z}{\mathbb{Z}}
\providecommand{\Q}{\mathbb{Q}}
\providecommand{\D}{\mathbb{D}}
\providecommand{\C}{\mathbb{C}}
\providecommand{\e}{{\mathrm{e}}}
\providecommand{\vect}[1]{\overrightarrow{#1}}
\providecommand{\vectt}[1]{\overrightarrow{\mathrm{#1}}}
\providecommand{\vvt}[1]{\overrightarrow{\mathrm{#1}}}
\providecommand{\np}[2][]{{#2~#1}}
\providecommand{\vv}[1]{{\overrightarrow{#1}}}
\renewcommand{\d}{{\mathrm{d}}}
\providecommand{\degres}{{^\circ}}
\providecommand{\degre}{{^\circ}}
\providecommand{\blue}{\color{blue}}
\providecommand{\red}{\color{red}}
\providecommand{\enskip}{\quad}
\providecommand{\llbracket}{\lbrack\!\lbrack}
\providecommand{\rrbracket}{\rbrack\!\rbrack}
\providecommand{\euro}{{\text{€}}}
\providecommand{\no}{{\text{n°}}}
\providecommand{\interval}[2]{{[#1;#2]}}
\providecommand{\intervFO}[2]{{[#1;#2[}}
\providecommand{\dotfill}{\ldots}

\providecommand{\ds}{\displaystyle}
\providecommand{\textding}[1]{#1}
\providecommand{\Oij}{\(\mathrm{O};\vec{\imath},\vec{\jmath}\)}
\providecommand{\Oijk}{\(\mathrm{O};\vec{\imath},\vec{\jmath},\vec{k}\)}
\providecommand{\Ouv}{\(\mathrm{O};\vec{u},\vec{v}\)}
\providecommand{\intervFF}[2]{\ensuremath{[#1;#2]}}
\renewcommand{\euro}{\ifmmode\text{\texteuro}\else\texteuro\fi}
\DeclareUnicodeCharacter{00B0}{\ifmmode^{\circ}\else\textdegree\fi}
\DeclareUnicodeCharacter{20AC}{\ifmmode\text{\texteuro}\else\texteuro\fi}
% Énoncés : Éducation nationale. Transcription et corrigés : APMEP, Denis Vergès et contributeurs. Découpage et classement : CoopMaths, Sébastien Lozano, Éric Elter et contributeurs. Adaptation de la présentation : Maths Explorer.
% CC BY-SA : https://coopmaths.fr/www/about/
% https://coopmaths.fr/alea/static/bac/2025/tex/bac_2025_05_sujet2_ameriquenord_1_cor.tex
\begin{document}
\noindent\textit{APMEP / CoopMaths. CC BY-SA. Présentation adaptée par Maths Explorer.}\par\medskip
\medskip



Les parties \textbf{A} et \textbf{B} sont indépendantes.

\bigskip

\textbf{Partie A}

\medskip






On note $N$ la variable aléatoire qui donne le nombre de paniers marqués.

\medskip

\emph{Les résultats des probabilités demandées seront, si nécessaire, arrondis au millième.}

\begin{enumerate}
\item 
L'épreuve est effectuée 15 fois de façon indépendante et à chaque lancer la probabilité de marquer est égale à $0,32$. $N$ suit donc la loi binomiale $\mathcal{B}(15~;~0,32)$.
\item 
La probabilité que Victor réussisse exactement 4 paniers est :

$P(N = 4) = \displaystyle \binom{15}{4}\times 0,32^4 \times (1 - 0,32)^{15-4} \approx 0,206$ au millième près.
\item 
La probabilité que Victor réussisse au plus 6 paniers est :

$P(N \leqslant 6) = P(N=0) + P(N = 1) + \ldots + P(N = 6) \approx 0,828$ (calculatrice).
\item 
On sait que : $E(N) = n \times p = 15 \times 0,32 = 4,8$.

Donc en moyenne sur 150 lancers, Victor marque 48 paniers.
\item 
	\begin{enumerate}
		\item 
Chaque lancer étant à 3 points, on a donc $T = 3N$.
		\item 
On a $E(T) = E (3N) = 3E(N)$\quad d'après la linéarité de l'espérance, soit

$E(T) = 3 \times 4,8 = 14,4$, donc une moyenne de 144 points marqués sur 150 lancers.
		\item 
On a $P(12 \leqslant T = 3N \leqslant 18) = P(4 \leqslant N \leqslant 6)$.

Cette probabilité est égale à :

$P(4 \leqslant N \leqslant 6) = P(N=4) + P(N=5) + P(N = 6) \approx 0,206 + 0,213 + 0,167$, soit

\[P(12 \leqslant T \leqslant 18) \approx 0,586.\]
	\end{enumerate}
\end{enumerate}

\medskip

\textbf{Partie B}

\medskip



















\begin{enumerate}
\item Dans cette question, on prend $n = 50$.

	\begin{enumerate}
		\item 
$M_{50}$ représente la moyenne empirique, sur 50 matchs, des points marqués par Victor
		\item 
$\bullet~$Chaque variable $X_1,\ldots, X_{50}$ suit la même loi que $X$ donc on a :

{$E(M_{50}) = E (X) = 22$ ;}

$\bullet~$ Comme les variables $X_1$, $X_2$,... $X_50$ sont indépendantes, on en déduit : \quad $V(M_{50}) = \dfrac{1}{50}V(X) = \dfrac{1}{50} \times 65 = \dfrac{65}{50} = \dfrac{130}{100} = 1,3$.

		\item 
D'après l'inégalité de Bienaymé-Tchébychev, pour une variable aléatoire $X$  d’espérance $E(X)$ et de variance $V (X)$,  pour tout réel $a > 0$ :
		$$P(|X - E(X)|\geqslant a) \leqslant \dfrac{V(X)}{a^2}$$

donc  ici  : $P(|M_{50} - E(M_{50})| \geqslant 3) \leqslant \dfrac{V(M_{50})}{3^2}$



Comme $\dfrac{V(M_{50})}{3^2} = \dfrac{1,3}{9} = \dfrac{13}{90}$, on a bien  :  $P\left(\left|M_{50}- 22\right| \geqslant 3\right) \leqslant \dfrac{13}{90}$.



		\item 
		L'événement \og  $19 <M_{50}< 25$ \fg{} revient à $\left|M_{50} - 22 \right| < 3$ et

$P\big(\left|M_{50} - 22 \right| < 3\big) = 1 - P\big(\left|M_{50} - 22 \right| \geqslant 3\big)$.

Soit $P\big(\left|M_{50} - 22 \right| \geqslant 3\big) \leqslant \dfrac{13}{90}$ entraîne $P\big(19 < M_{50} < 25 \big) \geqslant 1 - \dfrac{13}{90} = \dfrac{77}{90} \approx 0,856$ : cette probabilité est bien supérieure à 0,85.
	\end{enumerate}


\item 



D'après la loi faible des grands nombres, $\forall t \in \R^{*+}, \quad \lim\limits_{n \to +\infty} P\Big(\big|M_n - E(X)\big| \geqslant t\Big) = 0$, donc en particulier ici,  $\lim\limits_{n \to +\infty} P\Big(\big|M_n - 22\big| \geqslant 3\Big) = 0$ .


donc pour tout $r>0$ il existe un rang $N$ tel que pour tout $n \geqslant N$, $P\big(\left|M_{n}-22\right| \geqslant 3\big)< r$ et c'est vrai en particulier pour $r=0,01$.


L'affirmation est donc fausse.

\textbf{Autre méthode : }

Pour tout réel $t>0$, d'après l'inégalité de  concentration, on a :

$ P\big(\left|M_n-E(X)\right| \geqslant t\big) \leqslant \dfrac{V(X)}{t^2 n}$

	Pour $t=3$ on obtient :  $ P\big(|M_n-22| \geqslant 3\big)\leqslant \dfrac{65}{ n \times 3^2}$.



On veut donc que  :\quad $\begin{aligned}[t]
	 \dfrac{65}{9 n} \leqslant 0,01 &\iff
	 \dfrac{6500}{9} \leqslant n \\
	 &\iff 722+\dfrac{2}{9}\leqslant n
\end{aligned}$

L'inégalité est donc vraie pour tout entier $n \geqslant 723$.



\end{enumerate}



\bigskip

\end{document}
